\makeatletter
\def\input@path{{./}{../../../preamble/}{book2026-2027/preamble/}}
\makeatother
\input{preamble.tex}
\usepackage{needspace}
\hypersetup{pdftitle={Grade 10 Term 2 Discrete Probability}}
\setlength{\emergencystretch}{2em}
\title{Grade 10, Term 2, Lesson 2, Criterion C6\\Negative Binomial Trials to a Fixed Success Target}
\author{}
\date{}
\begin{document}
\maketitle
\setcounter{tocdepth}{1}
\setcounter{secnumdepth}{2}
\phantomsection\label{toc}
\tableofcontents
\section*{Evaluated criterion}
\textbf{C6:} Applies the negative binomial distribution to situations that count the number of trials needed to obtain a fixed number of successes.
\section{Introduction}\label{sec:t2-c6-introduction}
All data in this activity are hypothetical teaching models. Probabilities have no monetary units. Keep exact expressions where useful and round final probabilities to four decimal places.
\subsection{Wait for a target number of sales}
A school enterprise needs three sales. It contacts visitors independently, with constant purchase probability $p$. Let $T$ count all contacts up to and including the third sale. The target $r=3$ successes is fixed, but the number of trials is random. Support is $\{3,4,5,\ldots\}$ contacts for $0<p<1$.
\textbf{Check.} Could the third sale occur on contact two? \emph{Suggested answer:} No. One contact gives at most one success; $T\geq r$.


\subsection{Fix the final success}
For the third sale to occur on contact five, contact five must be a success. The first four contacts contain exactly two successes, in any positions. There are $\binom42=6$ choices of those positions, read as four choose two. In general
\[\binom ab=\frac{a!}{b!(a-b)!},\qquad4!=4\cdot3\cdot2\cdot1=24,\quad0!=1.\]
Factorial means multiply positive integers up to the number. Combinations count the selected success positions.
\textbf{Check.} Why not use $\binom53$? \emph{Suggested answer:} It would allow the third success before contact five; the final position must be fixed as success.


\subsection{Trials, failures and probability}
The failure count is $F=T-r$. If three sales require five contacts, there are two failures. Counts have units of contacts or failed contacts; probability has no units. A per-contact expense produces a separate monetary cost. Independence means earlier results do not change later probabilities.
\textbf{Check.} For $r=3$, $T=7$, what is $F$? \emph{Suggested answer:} $F=T-r=4$ failures.

\textbf{Check.} If $r=1$, which model results? \emph{Suggested answer:} Geometric, counting total trials until the first success.

\section{General worked problems}
\Needspace{14\baselineskip}
\subsection{Introductory general problem: second success on trial three}
Independent binary trials each succeed with probability $1/2$.

\textbf{Questions/Tasks}
\begin{enumerate}
\item Calculate $P(T=3)$. 
\end{enumerate}
\subsubsection*{Complete worked solution}
\begin{center}
\begin{tabular}{llp{0.48\linewidth}}
\toprule
Symbol & Value & Meaning in this problem \\ \midrule
$r$ & $2$ successes & Target number of successes \\
$p$ & $0.5$ & Success probability per trial \\
$t$ & $3$ trials & Requested stopping trial \\
\bottomrule
\end{tabular}
\end{center}

\par\noindent\begin{minipage}{\linewidth}
\textbf{Step 1: Define and identify.} $T$ is the total trials up to and including success 2. Support is $\{2,3,\ldots\}$ trials; $r=2$ and $p=0.5$.
\end{minipage}\par\medskip
\par\noindent\begin{minipage}{\linewidth}
\textbf{Step 2: Check the model.} Binary trials are independent with constant success probability and stop at the fixed target. This is the total-trial negative binomial convention.
\end{minipage}\par\medskip
\par\noindent\begin{minipage}{\linewidth}
\textbf{Step 3: Translate and calculate.} Use
\[P(T=t)=\binom{t-1}{r-1}p^r(1-p)^{t-r}.\]
Substitute:
\[P(T=3)=\binom{2}{1}(0.5)^2(0.5)^{1}\approx 0.2500.\]
\end{minipage}\par\medskip
\par\noindent\begin{minipage}{\linewidth}
\textbf{Step 4: Check and interpret.} The probability is in $[0,1]$; the earliest stopping trial is 2. On the exact event there are 1 failures, 1 successes in the first 2 trials, then the final success. Allowed sequences are SFS and FSS; SSF reaches the second success on trial 2.
\end{minipage}\par\medskip
\subsection{Bridge to the general model}
\[P(T=t)=\binom{t-1}{r-1}p^r(1-p)^{t-r},\quad t=r,r+1,\ldots.\]
The final trial is success. Choose $r-1$ success positions among the preceding $t-1$ trials; there are $t-r$ failures overall. Use this when independent binary trials with constant $0<p<1$ continue until success $r$. When $p=1$, $T=r$ certainly; when $p=0$, the target is never reached. For $r=1$, $\binom{t-1}{0}=1$, giving $(1-p)^{t-1}p$, the geometric PMF.

The sum $\sum_{u=r}^{m}$ adds exact stopping probabilities for $u=r,r+1,\ldots,m$. Alternatively $T\leq m$ means at least $r$ successes in the first $m$ trials; a binomial tail provides another calculation.
\Needspace{14\baselineskip}
\subsection{Advanced general problem: target by a deadline}
Independent trials have success probability 0.4. Stop at the third success.

\textbf{Questions/Tasks}
\begin{enumerate}
\item Calculate $P(T\leq6)$. Verify using a binomial count in the first six trials.
\end{enumerate}
\subsubsection*{Complete worked solution}
\par\noindent\begin{minipage}{\linewidth}
\textbf{Step 1: Define and identify.} $T$ is the total trials up to and including success 3. Support is $\{3,4,\ldots\}$ trials; $r=3$ and $p=0.4$.
\end{minipage}\par\medskip
\par\noindent\begin{minipage}{\linewidth}
\textbf{Step 2: Check the model.} Binary trials are independent with constant success probability and stop at the fixed target. This is the total-trial negative binomial convention.
\end{minipage}\par\medskip
\par\noindent\begin{minipage}{\linewidth}
\textbf{Step 3: Translate and calculate.} Use
\[P(T\leq m)=\sum_{u=r}^{m}\binom{u-1}{r-1}p^r(1-p)^{u-r}.\]
Substitute:
\[P(T\leq6)=\sum_{u=3}^{6}\binom{u-1}{2}(0.4)^3(0.6)^{u-3}\approx 0.4557.\]
\end{minipage}\par\medskip
\par\noindent\begin{minipage}{\linewidth}
\textbf{Step 4: Check and interpret.} The probability is in $[0,1]$; the earliest stopping trial is 3. The cumulative sum includes all stopping trials from 3 through 6. For $B\sim\operatorname{Bin}(6,0.4)$, $P(T\leq6)=P(B\geq3)=1-\sum_{j=0}^{2}\binom6j(0.4)^j(0.6)^{6-j}=0.4557$. Both describe reaching three successes by trial 6.
\end{minipage}\par\medskip
\section{Notation reference and practice procedure}
\begin{center}
\begin{tabular}{lp{0.72\linewidth}}
\toprule
Symbol & Meaning and units \\ \midrule
$r$ & Fixed target successes \\
$T,t$ & Total trials and one stopping trial \\
$p$ & Constant success probability per trial \\
$F=T-r$ & Failures before success $r$ \\
$\binom{t-1}{r-1}$ & Arrangements before the mandatory final success \\
\bottomrule
\end{tabular}
\end{center}
\begin{enumerate}
\item Define the random variable in words and give its units and support.
\item List parameters and check that the model assumptions match the situation.
\item Translate ordinary language into an event, paying attention to endpoints.
\item Write the symbolic formula first, then substitute without rounding intermediate values.
\item Check bounds, total probability or a complementary calculation; interpret the answer in context.
\end{enumerate}
\Needspace{19\baselineskip}
\section{Foundational contextual practice}\label{sec:t2-c6-foundational}
\Needspace{14\baselineskip}
\subsection{Two sales on call four}
Independent calls each sell with probability 0.3. A sale is success; stop at two sales.

\textbf{Questions/Tasks}
\begin{enumerate}
\item Calculate $P(T=4)$. 
\end{enumerate}
\subsubsection*{Complete worked solution}
\par\noindent\begin{minipage}{\linewidth}
\textbf{Step 1: Define and identify.} $T$ is the total trials up to and including success 2. Support is $\{2,3,\ldots\}$ trials; $r=2$ and $p=0.3$.
\end{minipage}\par\medskip
\par\noindent\begin{minipage}{\linewidth}
\textbf{Step 2: Check the model.} Binary trials are independent with constant success probability and stop at the fixed target. This is the total-trial negative binomial convention.
\end{minipage}\par\medskip
\par\noindent\begin{minipage}{\linewidth}
\textbf{Step 3: Translate and calculate.} Use
\[P(T=t)=\binom{t-1}{r-1}p^r(1-p)^{t-r}.\]
Substitute:
\[P(T=4)=\binom{3}{1}(0.3)^2(0.7)^{2}\approx 0.1323.\]
\end{minipage}\par\medskip
\par\noindent\begin{minipage}{\linewidth}
\textbf{Step 4: Check and interpret.} The probability is in $[0,1]$; the earliest stopping trial is 2. On the exact event there are 2 failures, 1 successes in the first 3 trials, then the final success. 
\end{minipage}\par\medskip
\secInicio[][sec:t2-c6-foundational]
\Needspace{14\baselineskip}
\subsection{Third approval on application five}
In a hypothetical independent stream each application has approval probability 0.6. Approval is success; stop at the third approval.

\textbf{Questions/Tasks}
\begin{enumerate}
\item Calculate $P(T=5)$. Find the rejection count on this event.
\end{enumerate}
\subsubsection*{Complete worked solution}
\par\noindent\begin{minipage}{\linewidth}
\textbf{Step 1: Define and identify.} $T$ is the total trials up to and including success 3. Support is $\{3,4,\ldots\}$ trials; $r=3$ and $p=0.6$.
\end{minipage}\par\medskip
\par\noindent\begin{minipage}{\linewidth}
\textbf{Step 2: Check the model.} Binary trials are independent with constant success probability and stop at the fixed target. This is the total-trial negative binomial convention.
\end{minipage}\par\medskip
\par\noindent\begin{minipage}{\linewidth}
\textbf{Step 3: Translate and calculate.} Use
\[P(T=t)=\binom{t-1}{r-1}p^r(1-p)^{t-r}.\]
Substitute:
\[P(T=5)=\binom{4}{2}(0.6)^3(0.4)^{2}\approx 0.2074.\]
\end{minipage}\par\medskip
\par\noindent\begin{minipage}{\linewidth}
\textbf{Step 4: Check and interpret.} The probability is in $[0,1]$; the earliest stopping trial is 3. On the exact event there are 2 failures, 2 successes in the first 4 trials, then the final success. $F=T-r=5-3=2$ rejected applications.
\end{minipage}\par\medskip
\secInicio[][sec:t2-c6-foundational]
\Needspace{19\baselineskip}
\section{Intermediate contextual practice}\label{sec:t2-c6-intermediate}
\Needspace{14\baselineskip}
\subsection{Second invoice settled by contact four}
Independent contacts each lead to a settlement with probability 0.5. Settlement is success; stop after two.

\textbf{Questions/Tasks}
\begin{enumerate}
\item Calculate $P(T\leq4)$. 
\end{enumerate}
\subsubsection*{Complete worked solution}
\par\noindent\begin{minipage}{\linewidth}
\textbf{Step 1: Define and identify.} $T$ is the total trials up to and including success 2. Support is $\{2,3,\ldots\}$ trials; $r=2$ and $p=0.5$.
\end{minipage}\par\medskip
\par\noindent\begin{minipage}{\linewidth}
\textbf{Step 2: Check the model.} Binary trials are independent with constant success probability and stop at the fixed target. This is the total-trial negative binomial convention.
\end{minipage}\par\medskip
\par\noindent\begin{minipage}{\linewidth}
\textbf{Step 3: Translate and calculate.} Use
\[P(T\leq m)=\sum_{u=r}^{m}\binom{u-1}{r-1}p^r(1-p)^{u-r}.\]
Substitute:
\[P(T\leq4)=\sum_{u=2}^{4}\binom{u-1}{1}(0.5)^2(0.5)^{u-2}\approx 0.6875.\]
\end{minipage}\par\medskip
\par\noindent\begin{minipage}{\linewidth}
\textbf{Step 4: Check and interpret.} The probability is in $[0,1]$; the earliest stopping trial is 2. The cumulative sum includes all stopping trials from 2 through 4. 
\end{minipage}\par\medskip
\secInicio[][sec:t2-c6-intermediate]
\Needspace{14\baselineskip}
\subsection{Below the insurance-sales target}
Independent offers each sell a policy with probability 0.2. Stop at three sales.

\textbf{Questions/Tasks}
\begin{enumerate}
\item Calculate $P(T>6)$. 
\end{enumerate}
\subsubsection*{Complete worked solution}
\par\noindent\begin{minipage}{\linewidth}
\textbf{Step 1: Define and identify.} $T$ is the total trials up to and including success 3. Support is $\{3,4,\ldots\}$ trials; $r=3$ and $p=0.2$.
\end{minipage}\par\medskip
\par\noindent\begin{minipage}{\linewidth}
\textbf{Step 2: Check the model.} Binary trials are independent with constant success probability and stop at the fixed target. This is the total-trial negative binomial convention.
\end{minipage}\par\medskip
\par\noindent\begin{minipage}{\linewidth}
\textbf{Step 3: Translate and calculate.} Use
\[P(T>m)=1-P(T\leq m).\]
Substitute:
\[P(T>6)=1-\left[\sum_{u=3}^{6}\binom{u-1}{2}(0.2)^3(0.8)^{u-3}\right]\approx 0.9011.\]
\end{minipage}\par\medskip
\par\noindent\begin{minipage}{\linewidth}
\textbf{Step 4: Check and interpret.} The probability is in $[0,1]$; the earliest stopping trial is 3. The cumulative sum includes all stopping trials from 3 through 6. The target has not been reached after six offers: fewer than three successes occurred in those offers.
\end{minipage}\par\medskip
\secInicio[][sec:t2-c6-intermediate]
\Needspace{19\baselineskip}
\section{Advanced contextual practice}\label{sec:t2-c6-advanced}
\Needspace{14\baselineskip}
\subsection{Sales within an outreach budget}
Each independent contact has sale probability 0.4 and costs 2 USD. The team needs three sales and has a 12 USD contact budget.

\textbf{Questions/Tasks}
\begin{enumerate}
\item Calculate $P(T\leq6)$. Translate the budget into a stopping-trial event.
\end{enumerate}
\subsubsection*{Complete worked solution}
\par\noindent\begin{minipage}{\linewidth}
\textbf{Step 1: Define and identify.} $T$ is the total trials up to and including success 3. Support is $\{3,4,\ldots\}$ trials; $r=3$ and $p=0.4$.
\end{minipage}\par\medskip
\par\noindent\begin{minipage}{\linewidth}
\textbf{Step 2: Check the model.} Binary trials are independent with constant success probability and stop at the fixed target. This is the total-trial negative binomial convention.
\end{minipage}\par\medskip
\par\noindent\begin{minipage}{\linewidth}
\textbf{Step 3: Translate and calculate.} Use
\[P(T\leq m)=\sum_{u=r}^{m}\binom{u-1}{r-1}p^r(1-p)^{u-r}.\]
Substitute:
\[P(T\leq6)=\sum_{u=3}^{6}\binom{u-1}{2}(0.4)^3(0.6)^{u-3}\approx 0.4557.\]
\end{minipage}\par\medskip
\par\noindent\begin{minipage}{\linewidth}
\textbf{Step 4: Check and interpret.} The probability is in $[0,1]$; the earliest stopping trial is 3. The cumulative sum includes all stopping trials from 3 through 6. $C=2T$ USD gives $C\leq12$ exactly when $T\leq6$. The fixed budget does not guarantee reaching the target.
\end{minipage}\par\medskip
\secInicio[][sec:t2-c6-advanced]
\Needspace{14\baselineskip}
\subsection{Failure counts and a geometric special case}
Independent recovery contacts succeed with probability 0.25. Stop at the second recovery; a report records three failures before it.

\textbf{Questions/Tasks}
\begin{enumerate}
\item Calculate $P(T=5)$. Translate three failures into total trials, then derive the model for target one.
\end{enumerate}
\subsubsection*{Complete worked solution}
\par\noindent\begin{minipage}{\linewidth}
\textbf{Step 1: Define and identify.} $T$ is the total trials up to and including success 2. Support is $\{2,3,\ldots\}$ trials; $r=2$ and $p=0.25$.
\end{minipage}\par\medskip
\par\noindent\begin{minipage}{\linewidth}
\textbf{Step 2: Check the model.} Binary trials are independent with constant success probability and stop at the fixed target. This is the total-trial negative binomial convention.
\end{minipage}\par\medskip
\par\noindent\begin{minipage}{\linewidth}
\textbf{Step 3: Translate and calculate.} Use
\[P(T=t)=\binom{t-1}{r-1}p^r(1-p)^{t-r}.\]
Substitute:
\[P(T=5)=\binom{4}{1}(0.25)^2(0.75)^{3}\approx 0.1055.\]
\end{minipage}\par\medskip
\par\noindent\begin{minipage}{\linewidth}
\textbf{Step 4: Check and interpret.} The probability is in $[0,1]$; the earliest stopping trial is 2. On the exact event there are 3 failures, 1 successes in the first 4 trials, then the final success. $F=3$, $r=2$ imply $T=F+r=5$. For $r=1$, the PMF becomes $\binom{t-1}{0}(0.25)(0.75)^{t-1}=(0.75)^{t-1}(0.25)$, the geometric convention.
\end{minipage}\par\medskip
\secInicio[][sec:t2-c6-advanced]
\section{Closing reference: key ideas and common errors}
The target $r$ is fixed and total trials $T$ are random. Include the final success. Use $\binom{t-1}{r-1}$, not $\binom tr$. Failures equal $T-r$. By-trial questions sum feasible stopping trials. Target one recovers geometric; fixed trials with random success count give binomial.
\secInicio
\end{document}
